% U4 WS2 -- net ionic equations, reaction types, acid-base.
\documentclass{apchem-worksheet}

\apcunit{4}
\apcunittitle{Chemical Reactions}
\apcdoctitle{Worksheet 2 \textbullet\ Net Ionic Equations and Reaction Types}
\apcsubtitle{CED 4.2, 4.7, 4.8 \textbullet\ decide strong vs weak before you write}

\begin{document}
\wsheader[State the reaction type first, then write the equation. For every
species you keep whole, be ready to say why.]

\problem[4.7]{Classify each reaction as precipitation, acid--base, or redox:}
\begin{parts}
  \item \ce{2 K(s) + Br2(l) -> 2 KBr(s)}:
        \answerblank[4cm]{redox}
  \item \ce{HNO3(aq) + LiOH(aq) -> LiNO3(aq) + H2O(l)}:
        \answerblank[4cm]{acid--base}
  \item \ce{Pb(NO3)2(aq) + K2CrO4(aq) -> PbCrO4(s) + 2 KNO3(aq)}:
        \answerblank[4cm]{precipitation}
\end{parts}

\problem[4.2]{Write the complete ionic and net ionic equations for
\ce{Ba(NO3)2(aq) + K2SO4(aq) -> BaSO4(s) + 2 KNO3(aq)}:}
\begin{parts}
  \item Complete ionic: \answerlines[2]{\ce{Ba^2+(aq) + 2 NO3^-(aq) +
        2 K+(aq) + SO4^2-(aq) -> BaSO4(s) + 2 K+(aq) + 2 NO3^-(aq)}}
  \item Spectator ions: \answerblank[4.5cm]{\ce{K+} and \ce{NO3-}}
  \item Net ionic: \answerblank[6cm]{\ce{Ba^2+ + SO4^2- -> BaSO4(s)}}
\end{parts}

\problem[4.2]{Write net ionic equations:}
\begin{parts}
  \item \ce{Na2S(aq) + Zn(NO3)2(aq)}:
        \answerblank[6.5cm]{\ce{Zn^2+ + S^2- -> ZnS(s)}}
  \item \ce{HBr(aq) + NaOH(aq)}:
        \answerblank[6.5cm]{\ce{H+ + OH^- -> H2O(l)}}
  \item \ce{HC2H3O2(aq) + KOH(aq)}:
        \answerlines[2]{\ce{HC2H3O2(aq) + OH^-(aq) -> C2H3O2^-(aq) +
        H2O(l)}}
  \item \ce{NH3(aq) + HCl(aq)}:
        \answerblank[6.5cm]{\ce{NH3 + H+ -> NH4+}}
\end{parts}

\problem[4.2]{Explain why (b) and (c) above have different net ionic
equations, even though both are ``an acid reacting with a strong base.''
\answerlines[3]{\ce{HBr} is a strong acid, fully ionized, so the reacting
species is free \ce{H+}. Acetic acid is weak and remains almost entirely
molecular, so the whole molecule must appear and transfer its proton
directly.}}

\problem[4.7]{Predict the products, write the balanced formula equation, and
state whether a precipitate forms:}
\begin{parts}
  \item \ce{AgNO3(aq) + K2CrO4(aq)}
        \answerlines[2]{\ce{2 AgNO3(aq) + K2CrO4(aq) -> Ag2CrO4(s) +
        2 KNO3(aq)}; yes, silver chromate is insoluble.}
  \item \ce{NaCl(aq) + Mg(NO3)2(aq)}
        \answerlines[2]{No reaction --- \ce{NaNO3} and \ce{MgCl2} are both
        soluble, so no precipitate forms.}
\end{parts}

\problem[4.8]{Classify and explain:}
\begin{parts}
  \item Is \ce{NH3} an Arrhenius base? A Bronsted--Lowry base? Explain the
        difference. \answerlines[3]{Not straightforwardly Arrhenius --- it
        contains no \ce{OH-}. It is clearly Bronsted--Lowry, because it
        accepts a proton to form \ce{NH4+}. (It does raise \ce{OH-} in
        water indirectly, which is why the broader definition is more
        useful.)}
  \item Name the six strong acids.
        \answerlines[2]{\ce{HCl}, \ce{HBr}, \ce{HI}, \ce{HNO3},
        \ce{H2SO4}, \ce{HClO4}}
\end{parts}

\problem[4.2]{Three different reactant pairs all give the net ionic equation
\ce{Ag+(aq) + Cl^-(aq) -> AgCl(s)}. Write formula equations for two of them.
\answerlines[3]{e.g.\ \ce{AgNO3(aq) + NaCl(aq) -> AgCl(s) + NaNO3(aq)} and
\ce{AgNO3(aq) + KCl(aq) -> AgCl(s) + KNO3(aq)}.}}

\begin{spiralstrip}{Unit 4 \textbullet\ blocks 1--2}
  \item Physical or chemical: dry ice subliming?
        \answerblank[2.4cm]{physical}
  \item Conserve what two things in a particulate drawing?
        \answerblank[3.6cm]{atoms and charge}
  \item Soluble? \ce{PbCl2} \answerblank[2.4cm]{no}
  \item Always-soluble cation family: \answerblank[4.2cm]{alkali metals,
        \ce{NH4+}}
  \item Is dissolving \ce{HCl(g)} in water physical or chemical?
        \answerblank[4.2cm]{chemical (it ionizes)}
\end{spiralstrip}

\end{document}
